\scanpage{12}
\mspage{P 49}
celle du pionnier ou de l'explorateur, parti à la découverte et à l'exploration de mondes inconnus et sans nom, m'appelant sans cesse pour que je les connaisse et les nomme. Cette pulsion-là, et l'énergie que j'y investissais (comme à la dérobée, quasiment !), étaient constamment à la portion congrue.

Pourtant, je savais bien au fond que c'était cette énergie-là, dérobée (pour ainsi dire) à celle que je devais à mes ``tâches'', qui était de l'essence la plus rare et la plus déliée - que la ``création'' dans mon travail de mathématicien, c'était avant tout l à qu'elle se plaçait : dans cette attention intense pour appréhender, dans les replis obscurs, informes et moites d'une chaude et inépuisable matrice nourricière, les premières traces de forme et de contours de ce qui n'était pas né encore et qui semblait m'appeler, pour prendre forme et s'incarner et naître... Dans le travail de découverte, cette attention intense, cette sollicitude ardente sont une force essentielle, tout comme la chaleur du soleil pour l'obscure gestation des semences enfouies dans la terre nourricière, et pour leur humble et miraculeuse éclosion à la lumière du jour.

Dans mon travail de mathématicien, je vois à l'oeuvre surtout ces deux forces ou pulsions, également profondes, de nature (me semble\unsure{-{}-|-}t-il) différentes. Pour évoquer l'une et l'autre, j'ai utilisé l'image du b â t i s s e u r , et celle du p i o n n i e r ou de l'explorateur. Mises côte-à-côte, l'une et l'autre me frappent soudain comme vraiment très ``yang'', très ``masculines'', voire ``macho'' ! Elles ont la résonance altière du mythe, ou celle des ``grandes occasions''. Sûrement elles sont inspirées par les vestiges, en moi, de mon ancienne vision ``héroïque'' du travail créateur, la vision super-yang. Telles quelles, elles donnent une vision fortement\unsure{·} teintée, pour ne pas dire figée, ``au garde à vous'', d'une réalité bien plus fluide, plus humble, plus ``simple'' - d'une réalité v i v a n t e .

Dans cette mâle pulsion du ``bâtisseur'', qui semble sans cesse me pousser vers de nouveaux chantiers, je discerne bien pourtant, en même temps, celle du c a s a n i e r : de celui profondément attaché à `` l a '' maison. Avant toute autre chose, c'est `` s a '' maison, celle des `` p r o c h e s '' - le lieu d'une intime entité vivante dont il se sent faire partie. Ensuite seulement, et à mesure que s'élargit le cercle de ce qui est ressenti comme ``proche'', est-elle aussi une ``maison pour tous''. Et dans cette pulsion de ``faire des maisons'' (comme on ``ferait'' l'amour...) il y a aussi et avant tout une t e n d r e s s e. Il y a la pulsion du c o n t a c t avec ces matériaux qu'on façonne un à un, avec un soin amoureux, et qu'on ne connaît vraiment que par ce contact aimant. Et, une fois montés les murs et posés les poutres et le toit, il y a la satisfaction
% archive sheet number: 11

